\chapter{Reference tables and examination review}

\section{How to use the tables}
The tables in this chapter use approximate United States reference values for generally healthy adults aged 19--50. Values differ by sex, age, pregnancy, lactation, country, and guideline. They are for education, not individualized prescribing. The UL applies only when established and may refer to supplements rather than food. Always check the current national reference source before clinical use. \cite{nap}

\section{Adult vitamin reference values}
\small
\begin{longtable}{p{0.18\textwidth}p{0.18\textwidth}p{0.18\textwidth}p{0.34\textwidth}}
\caption{Approximate adult vitamin reference values}\\
\toprule
Nutrient & Adult male & Adult female & Principal clinical note \\
\midrule
\endfirsthead
\toprule Nutrient & Adult male & Adult female & Principal clinical note \\
\midrule
\endhead
Vitamin A & 900 mcg RAE & 700 mcg RAE & Vision, epithelium, immunity; preformed retinol has a UL.\\
Thiamin (B1) & 1.2 mg & 1.1 mg & Oxidative metabolism; deficiency may cause Wernicke disease.\\
Riboflavin (B2) & 1.3 mg & 1.1 mg & FMN/FAD cofactor; deficiency often accompanies malnutrition.\\
Niacin (B3) & 16 mg NE & 14 mg NE & NAD/NADP; pharmacological doses require supervision.\\
Pantothenic acid (B5) & 5 mg AI & 5 mg AI & Coenzyme A; deficiency is rare.\\
Vitamin B6 & 1.3 mg & 1.3 mg & Amino-acid and neurotransmitter metabolism; excess can cause neuropathy.\\
Biotin (B7) & 30 mcg AI & 30 mcg AI & Carboxylase cofactor; excess can interfere with tests.\\
Folate (B9) & 400 mcg DFE & 400 mcg DFE & DNA synthesis; folic acid is important before conception.\\
B12 & 2.4 mcg & 2.4 mcg & Hematology and nervous system; absorption is complex.\\
Vitamin C & 90 mg & 75 mg & Collagen and iron absorption; smokers need additional intake.\\
Vitamin D & 15 mcg & 15 mcg & Calcium-phosphate regulation; 25(OH)D is the usual status marker.\\
Vitamin E & 15 mg & 15 mg & Membrane antioxidant; deficiency is uncommon.\\
Vitamin K & 120 mcg AI & 90 mcg AI & Coagulation; consistency is important with warfarin.\\
\bottomrule
\end{longtable}
\normalsize

\section{Adult mineral reference values}
\small
\begin{longtable}{p{0.20\textwidth}p{0.18\textwidth}p{0.18\textwidth}p{0.32\textwidth}}
\caption{Approximate adult mineral reference values}\\
\toprule
Nutrient & Adult male & Adult female & Principal clinical note \\
\midrule
\endfirsthead
\toprule Nutrient & Adult male & Adult female & Principal clinical note \\
\midrule
\endhead
Calcium & 1,000 mg & 1,000 mg & Bone, muscle, nerve, clotting; kidney disease changes management.\\
Phosphorus & 700 mg & 700 mg & Bone, ATP, membranes; excess is important in kidney disease.\\
Magnesium & 400 mg & 310 mg & Enzyme cofactor; supplemental UL is lower than total dietary need.\\
Iron & 8 mg & 18 mg & Hemoglobin; investigate blood loss in adult iron-deficiency anemia.\\
Zinc & 11 mg & 8 mg & Immunity and wound healing; excess zinc can cause copper deficiency.\\
Copper & 900 mcg & 900 mcg & Iron metabolism and nervous system; assess with context.\\
Iodine & 150 mcg & 150 mcg & Thyroid hormone; both deficiency and excess matter.\\
Selenium & 55 mcg & 55 mcg & Selenoproteins; Brazil nuts may be highly concentrated.\\
Manganese & 2.3 mg AI & 1.8 mg AI & Enzyme systems; deficiency is uncommon.\\
Chromium & 35 mcg AI & 25 mcg AI & Metabolism; no routine supplement indication.\\
Molybdenum & 45 mcg & 45 mcg & Enzyme cofactor; deficiency is rare.\\
Fluoride & 4 mg AI & 3 mg AI & Dental and bone mineralization; excess causes fluorosis.\\
Potassium & 3,400 mg AI & 2,600 mg AI & Cardiac and neuromuscular physiology; kidney disease changes safety.\\
\bottomrule
\end{longtable}
\normalsize

\section{Clinical summary table}
\begin{longtable}{p{0.15\textwidth}p{0.24\textwidth}p{0.22\textwidth}p{0.20\textwidth}}
\caption{Selected deficiency and toxicity patterns}\\
\toprule
Nutrient & Deficiency pattern & Excess pattern & First clinical action \\
\midrule
\endfirsthead
\toprule Nutrient & Deficiency pattern & Excess pattern & First clinical action \\
\midrule
\endhead
A & Night blindness, xerophthalmia & Liver, bone, fetal toxicity & Review retinoid exposure and pregnancy.\\
D & Rickets, osteomalacia & Hypercalcemia & Check calcium, kidney function, and dose.\\
B1 & Wernicke, neuropathy, heart failure & Usually low toxicity concern & Treat urgently when suspected; investigate cause.\\
B6 & Dermatitis, anemia, neuropathy & Sensory neuropathy & Stop unnecessary high-dose products.\\
B9/B12 & Megaloblastic anemia & Folate may obscure B12 disease & Test and treat the correct deficiency.\\
C & Scurvy & GI effects, stone risk in susceptible people & Assess diet and bleeding/wound signs.\\
Iron & Microcytic anemia, fatigue & GI and organ toxicity & Confirm iron deficiency and source of loss.\\
Calcium & Bone loss, hypocalcemia & Stones, hypercalcemia & Review intake, kidney disease, medicines.\\
Iodine & Goiter, thyroid dysfunction & Thyroid dysfunction & Review supplements and thyroid tests.\\
\bottomrule
\end{longtable}

\section{High-yield examination points}
\begin{enumerate}
\item Fat-soluble vitamins are more likely to accumulate, but water-soluble vitamins are not universally safe at high doses.
\item Iron-deficiency anemia in an adult requires evaluation for blood loss, not only dietary advice.
\item Neurologic B12 deficiency can occur without anemia.
\item Vitamin K intake should generally be consistent during warfarin therapy.
\item Potassium and magnesium supplements can be hazardous in kidney disease.
\item Folic acid prevents neural-tube defects; it does not treat every cause of macrocytosis.
\item A normal serum nutrient value does not always exclude tissue deficiency or functional disease.
\end{enumerate}

\section{Sample multiple-choice questions}
\textbf{1.} A patient with alcohol use disorder presents with confusion, ataxia, and ophthalmoplegia. Which nutrient deficiency is most urgent to consider?\\
\textit{Answer: Thiamin deficiency causing Wernicke encephalopathy.}

\textbf{2.} Which test best reflects vitamin D stores in ordinary clinical assessment?\\
\textit{Answer: Serum 25-hydroxyvitamin D, interpreted with clinical context.}

\textbf{3.} A vegan patient has paresthesias with a normal hemoglobin. Which deficiency remains important?\\
\textit{Answer: Vitamin B12 deficiency.}

\textbf{4.} A patient taking warfarin wants to eliminate all green vegetables. What is the best advice?\\
\textit{Answer: Maintain a consistent vitamin K intake and coordinate changes with the anticoagulation team.}

\textbf{5.} Which supplement is particularly associated with interference in some immunoassays?\\
\textit{Answer: High-dose biotin.}

\section{Short-answer prompts}
Explain why serum calcium does not directly measure bone calcium; compare food folate with folic acid; describe three causes of iron deficiency; explain how kidney disease changes mineral safety; and outline a safe supplement review for a pregnant patient.

\section{Answer-writing framework}
Strong examination answers define the nutrient, name its active form or major pathway, connect the pathway to organ function, describe deficiency and excess, identify high-risk groups, and finish with assessment and management principles. Avoid presenting a supplement dose without specifying age, indication, formulation, contraindications, and monitoring.
